\chapter{Étape 3 : Architecture du Système}

\section{Vue d'Ensemble}

Le système GRM adopte une architecture \textbf{3-tiers} classique, séparant clairement :
\begin{itemize}
    \item Le \textbf{Frontend} (couche présentation) : application React SPA
    \item Le \textbf{Backend} (couche logique métier) : API REST Node.js/Express
    \item La \textbf{Base de données} (couche persistance) : PostgreSQL
\end{itemize}

Cette architecture garantit la séparation des responsabilités, la maintenabilité et la scalabilité.

\section{Architecture 3-Tiers}

\begin{verbatim}
+----------------------------------------------------------+
|                     CLIENT (Browser)                     |
|                  React SPA (Port 3000)                   |
|  +------------------+   +-----------------------------+  |
|  | Pages / Composants|   |  Context + Services (Axios) | |
|  +------------------+   +-----------------------------+  |
+------------------------------|---------------------------+
                               | HTTPS / REST API
+------------------------------|---------------------------+
|                    BACKEND (Node.js)                     |
|                  Express API (Port 5000)                  |
|  +----------+  +-------------+  +-------------------+   |
|  | Routes   |--| Controllers |--| Middleware (JWT)   |   |
|  +----------+  +-------------+  +-------------------+   |
+------------------------------|---------------------------+
                               | SQL (pg)
+------------------------------|---------------------------+
|               BASE DE DONNÉES (PostgreSQL)               |
|                       Port 5432                          |
|  +----------+  +----------+  +----------+  +--------+   |
|  | users    |  | tenders  |  | resources|  | offers |   |
|  +----------+  +----------+  +----------+  +--------+   |
+----------------------------------------------------------+
\end{verbatim}

\section{Choix Technologiques}

\begin{table}[h!]
\centering
\begin{tabularx}{\textwidth}{|l|l|X|}
\hline
\textbf{Couche} & \textbf{Technologie} & \textbf{Justification} \\
\hline
Frontend & React 18 & Bibliothèque UI déclarative, SPA, grande communauté, React Router v6 pour la navigation \\
\hline
HTTP Client & Axios & Client HTTP avec intercepteurs pour gestion centralisée des tokens JWT \\
\hline
Backend & Node.js + Express & Performance, non-bloquant, JSON natif, facilité d'API REST \\
\hline
Authentification & JWT (jsonwebtoken) & Stateless, sécurisé, contient le rôle et l'ID pour le RBAC \\
\hline
Hachage & bcrypt & Standard industrie pour le hachage sécurisé des mots de passe \\
\hline
Base de données & PostgreSQL & SGBDR robuste, JSONB pour les specs flexibles, UUID pour les clés primaires \\
\hline
Driver DB & node-postgres (pg) & Accès direct SQL avec pool de connexions \\
\hline
Déploiement & Docker & Containerisation pour reproductibilité des environnements \\
\hline
\end{tabularx}
\caption{Stack technologique}
\end{table}

\section{Diagramme de Composants}

\begin{verbatim}
React Frontend
+--------------------------------------------+
|  <<component>>         <<component>>       |
|  AuthContext    <-->   API Service (Axios)  |
|                              |              |
|  <<component>>   <<component>>  <<page>>   |
|  Layout       | TenderList  | Dashboard    |
|  (Navbar,     | TenderDetail| ResourceList |
|   Sidebar)    | TenderForm  | BreakdownList|
+--------------------------------------------+
                       |
              REST API (JSON/HTTP)
                       |
Express Backend
+--------------------------------------------+
|  server.js                                 |
|     |                                      |
|  <<middleware>>  <<middleware>>             |
|  CORS          JWT authenticate            |
|  JSON Parser   RBAC authorize              |
|                                            |
|  <<routes>>  ----->  <<controllers>>       |
|  /api/auth           authController        |
|  /api/tenders        tenderController      |
|  /api/offers         offerController       |
|  /api/resources      resourceController    |
|  /api/suppliers      supplierController    |
|  /api/breakdowns     breakdownController   |
|  /api/departments    departmentController  |
+--------------------------------------------+
                       |
              node-postgres Pool
                       |
PostgreSQL Database
+--------------------------------------------+
|  users | departments | suppliers           |
|  tenders | tender_needs | offers           |
|  offer_items | resources | assignments     |
|  breakdowns | maintenance_reports          |
+--------------------------------------------+
\end{verbatim}

\section{Sécurité}

\subsection{Authentification JWT}

\begin{enumerate}
    \item L'utilisateur envoie email + mot de passe à \texttt{POST /api/auth/login}
    \item Le serveur vérifie le mot de passe avec \texttt{bcrypt.compare()}
    \item En cas de succès, le serveur génère un JWT signé (durée : 24h) contenant :
    \begin{itemize}
        \item \texttt{id} : identifiant utilisateur
        \item \texttt{role} : rôle pour le RBAC
        \item \texttt{department\_id} : département associé
        \item \texttt{supplier\_id} : ID fournisseur (si applicable)
    \end{itemize}
    \item Le client stocke le token dans \texttt{localStorage} et l'envoie dans l'en-tête \texttt{Authorization: Bearer <token>}
    \item Le middleware \texttt{authenticate} vérifie le token à chaque requête protégée
\end{enumerate}

\subsection{Contrôle d'Accès par Rôle (RBAC)}

\begin{table}[h!]
\centering
\begin{tabularx}{\textwidth}{|X|l|l|l|l|l|}
\hline
\textbf{Endpoint} & \textbf{Resp.} & \textbf{Chef} & \textbf{Ensgt} & \textbf{Tech.} & \textbf{Fsseur} \\
\hline
GET /tenders & \checkmark & \checkmark & & & \checkmark \\
\hline
POST /tenders & \checkmark & \checkmark & & & \\
\hline
PATCH /tenders/:id/close & \checkmark & & & & \\
\hline
POST /offers/tender/:id & & & & & \checkmark \\
\hline
PATCH /offers/:id/accept & \checkmark & & & & \\
\hline
POST /resources & \checkmark & & & & \\
\hline
POST /resources/:id/assign & \checkmark & & & & \\
\hline
PATCH /suppliers/:id/blacklist & \checkmark & & & & \\
\hline
POST /breakdowns & & \checkmark & \checkmark & & \\
\hline
POST /breakdowns/:id/maintenance-report & & & & \checkmark & \\
\hline
PATCH /breakdowns/:id/return-to-supplier & \checkmark & & & & \\
\hline
\end{tabularx}
\caption{Matrice des autorisations}
\end{table}

\section{Diagramme de Déploiement}

\begin{verbatim}
+---------------------+        +----------------------+
|   Machine Client    |        |  Serveur Applicatif  |
|  (Navigateur Web)   |        |  (Node.js / Docker)  |
|                     |        |                      |
|  React SPA          | HTTPS  |  Express API         |
|  (Build statique)   |------->|  Port 5000           |
|                     |        |                      |
+---------------------+        +----------+-----------+
                                          |
                                          | TCP/5432
                                          |
                               +----------+-----------+
                               |  Serveur PostgreSQL  |
                               |  (Docker / VM)       |
                               |  Base: grm_db        |
                               +----------------------+
\end{verbatim}

\subsection{Configuration Docker}

Le déploiement utilise Docker Compose avec trois services :
\begin{enumerate}
    \item \textbf{db} : PostgreSQL 15, volume persistant pour les données
    \item \textbf{backend} : Node.js 20-alpine, expose le port 5000
    \item \textbf{frontend} : Nginx servant le build React, expose le port 80
\end{enumerate}

\section{Diagramme de Séquence — Sélection d'une Offre}

\begin{verbatim}
Responsable    Frontend       Backend API      Database
     |             |               |               |
     |--[Voir offres d'un AO]----->|               |
     |             |   GET /offers/tender/{id}     |
     |             |               |--[SELECT offers JOIN suppliers]-->|
     |             |               |<--------[Liste offres triées]-----|
     |             |<-[Affiche offres par prix]    |               |
     |             |               |               |               |
     |--[Clic Accepter offre X]--->|               |               |
     |             |   PATCH /offers/{id}/accept   |               |
     |             |               |--[BEGIN]----->|               |
     |             |               |--[UPDATE offers SET status='accepted' WHERE id=X]-->|
     |             |               |--[UPDATE offers SET status='rejected' WHERE tender_id=T AND id!=X]-->|
     |             |               |--[UPDATE tenders SET status='awarded']-->|
     |             |               |--[COMMIT]---->|               |
     |             |<-[200 OK Offre acceptée]      |               |
     |<--[Notif: Offre acceptée/rejetée]           |               |
\end{verbatim}
